\emph{Aan alle idioten, die zijn of nog zullen komen, TABÉ.}

\emph{De volledige Volkvergadering heeft unaniem en met één stem
besloten:}

\begin{longtable}[]{@{}l@{}}
\toprule
\endhead
\begin{minipage}[t]{0.36\columnwidth}\raggedright
Als titel V van de Grondwet worden volgende artikelen aangenomen.\strut
\end{minipage}\tabularnewline
\begin{minipage}[t]{0.36\columnwidth}\raggedright
\#\#\#\# Art. 51 \emph{Goedgekeurd op 10/11/2025} Een groet voor,
tijdens of juist na een fysieke vergadering zal een handdruk zijn met de
linkerhand op even dagen en met de rechterhand op oneven dagen. De
Linker groet altijd met de linkerhand.\strut
\end{minipage}\tabularnewline
\begin{minipage}[t]{0.36\columnwidth}\raggedright
\#\#\#\# Art. 52 \emph{Goedgekeurd op 5/08/2025} Het Rijkswapen zal de
Opposum met een cowboyhoed zijn. De kenspreuk zal ``Be rootin, be
tootin, and by god be shooting but most of all be kind'' zijn.\strut
\end{minipage}\tabularnewline
\begin{minipage}[t]{0.36\columnwidth}\raggedright
\#\#\#\# Art. 53 Het Volkslied zal ``Bij de Preikswacht'' zijn. Dit lied
zal geschreven worden door Buriku Luigi Mussolini. Indien een Idioot het
Volkslied meezingt tijdens een nationale feestdag zal hij met een prei
moeten zwaaien.\strut
\end{minipage}\tabularnewline
\begin{minipage}[t]{0.36\columnwidth}\raggedright
\#\#\#\# Art. 54 De nationale klederdracht zal ten minste bestaan uit
een Hawaii hemd.\strut
\end{minipage}\tabularnewline
\begin{minipage}[t]{0.36\columnwidth}\raggedright
De leden der kamers voegen 4 gulden kwastjes toe aan elke mouw van hun
hemd. Buriku de heer Snetjes koopt deze aan voor iedereen. Buriku Luigi
Mussolini zal deze aan de hemden naaien.\strut
\end{minipage}\tabularnewline
\begin{minipage}[t]{0.36\columnwidth}\raggedright
De mandaat pin en eretekens worden op de linkerborst gedragen.\strut
\end{minipage}\tabularnewline
\begin{minipage}[t]{0.36\columnwidth}\raggedright
\#\#\#\# Art. 54bis Naast een hemd dragen sommige mandaathouders en
dignitarissen nog extra regalia met hun nationale klederdracht.\strut
\end{minipage}\tabularnewline
\begin{minipage}[t]{0.36\columnwidth}\raggedright
De Voorzitter van de Volksvergadering en de Voorzitter van de Andere
Volksvergadering dragen elk een Stok der Sprekers om het lid aan te
duiden die op dat moment het woord heeft. \strut
\end{minipage}\tabularnewline
\begin{minipage}[t]{0.36\columnwidth}\raggedright
De Linker draagt het Wetboek der Idioten op papier.\strut
\end{minipage}\tabularnewline
\begin{minipage}[t]{0.36\columnwidth}\raggedright
De Eerste Minister en de Minister van Alles dragen elk een zak met
knikkers omdat ze toch al met onze ballen spelen. \strut
\end{minipage}\tabularnewline
\begin{minipage}[t]{0.36\columnwidth}\raggedright
De Pontifex Maximus zal een eikenhouten kruis rond zijn nek dragen. Hoe
groter het kruis hoe heiliger de Pontifex Maximus zal zijn.\strut
\end{minipage}\tabularnewline
\begin{minipage}[t]{0.36\columnwidth}\raggedright
\#\#\#\# Art. 55 \emph{Goedgekeurd op 5/08} De nationale feestdagen
zullen 1 april en 11 september zijn.\strut
\end{minipage}\tabularnewline
\begin{minipage}[t]{0.36\columnwidth}\raggedright
\#\#\#\# Art. 56 Het nationale gerecht zal preischotel zonder prei
zijn.\strut
\end{minipage}\tabularnewline
\bottomrule
\end{longtable}

\emph{De Kontsuls van het Jaar des Herens Tweeduizendzesentwintig}

\emph{Op den \texttt{\textless{}Datum\textgreater{}}}
